% Generated draft: requires mathematical and source review.
\documentclass[11pt]{article}
\usepackage[margin=25mm]{geometry}
\usepackage{fontspec}
\setmainfont{Latin Modern Roman}
\usepackage{amsmath,amssymb,mathtools,unicode-math}
\setmathfont{Latin Modern Math}
\usepackage{enumitem,needspace,xurl,hyperref}
\setlength{\parindent}{0pt}
\setlength{\parskip}{5pt}
\setlength{\emergencystretch}{3em}
\newcommand{\N}{\mathbb{N}}
\newcommand{\Z}{\mathbb{Z}}
\newcommand{\Q}{\mathbb{Q}}
\newcommand{\R}{\mathbb{R}}
\newcommand{\C}{\mathbb{C}}
\newcommand{\F}{\mathbb{F}}
\newcommand{\A}{\mathbb{A}}
\newcommand{\PP}{\mathbb{P}}
\newcommand{\E}{\mathbb{E}}
\newcommand{\eps}{\varepsilon}
\newcommand{\dd}{\,\mathrm{d}}
\begin{document}
\section*{UnsolvedMath ID: 30005107}
\textbf{Source problem code:} catalogue entry
\subsection{Definitions and mathematical statement}
Let \(n\geq 1\) be an integer, and let \(K_n\) be the complete simple undirected graph with vertex set \([n]=\{1,\ldots,n\}\). A tree means a finite simple undirected connected graph with no cycles; in particular, the one-vertex graph is a tree. An edge decomposition into trees permits their embedded vertex sets to overlap, but requires their edge sets to partition \(E(K_n)\).

The conjecture attributed in the supplied report to Gy\'arf\'as is the following: for every integer \(n\geq 1\) and every sequence of trees \(T_1,\ldots,T_n\) satisfying \(|V(T_i)|=i\) for each \(i\in\{1,\ldots,n\}\), there exist injective maps
\[
\varphi_i:V(T_i)\longrightarrow[n]\qquad(1\leq i\leq n)
\]
such that, on defining
\[
F_i=\bigl\{\{\varphi_i(u),\varphi_i(v)\}:\{u,v\}\in E(T_i)\bigr\},
\]
one has
\[
E(K_n)=\bigcup_{i=1}^{n}F_i,
\qquad
F_i\cap F_j=\varnothing\quad\text{whenever }i\ne j.
\]
The copy of \(T_1\) contributes no edges, and for \(n=1\) the required edge partition is empty. Copies are subgraphs, not necessarily induced subgraphs.

The mathematical target is specified by the supplied passage, but its original reference labeled [6] is not bibliographically identified in the record. The supplied open status and verification labels are curation claims, not independent verification in this draft.
\subsection{Short English statement}
For every positive integer \(n\), can every choice of trees having respectively \(1,2,\ldots,n\) vertices be embedded in \(K_n\) so that each edge belongs to exactly one embedded tree? The trees may share vertices. The supplied record attributes this conjecture to Gy\'arf\'as and labels it open.
\subsection{Sources}
\begin{itemize}
\item Local UnsolvedMath registry, permanent ID \(30005107\), source problem code \texttt{catalogue entry}, title ``Gy\'arf\'as Tree-Packing Conjecture,'' in the Oberwolfach Reports Open Problems collection. The supplied registry record labels the statement \texttt{corrected\_verified} and says that the official report was checked directly; that verification is reported here as supplied metadata, not as a check performed for this draft.
\item Supplied report citation: \emph{Combinatorics, Probability and Computing (2023)}, Oberwolfach Reports, pp.~1165--1237, DOI \texttt{10.4171/owr/2022/22}. \url{https://doi.org/10.4171/owr/2022/22}. The excerpt attributes the conjecture to Gy\'arf\'as through reference [6], whose bibliographic entry is not supplied. Its discussion of Ringel's conjecture is historical context, not a separately verified current status claim.
\item Supplied literature source: Bottcher et al., \emph{Journal of Combinatorial Theory B}. \url{https://doi.org/10.1016/j.jctb.2016.03.003}.
\item Supplied literature source: Montgomery, Pokrovskiy, and Sudakov. \url{https://doi.org/10.1007/s00039-021-00576-2}.
\item Supplied literature assessment dated exactly \texttt{2026-08-22}: the arbitrary-sequence conjecture remains open, with bounded-degree and other substantial special cases known. This is supplied dated curation derived from the OWR record, not an independently verified literature assessment; its status and cited claims require independent checking.
\end{itemize}
\end{document}
